\section{The 3D-SCPTOP formulation}
\label{sec:model}

The model is written over \emph{trucks}. A dispatched truck designs its
own layout --- the stack types, orientations, positions and layer
compositions of its slots are variables of that truck and that period,
and its load is assigned to those slots in the same period --- so there
is no catalogue of loading configurations. That is the decision
Section~\ref{sec:decomp} rests on: every row below carries either the
triple $(g,k,t)$ or none of it, so the model is block-angular over
$G\times K\times T$ and the pricing problem of the Dantzig--Wolfe
reformulation is obtained by deleting that quantifier and nothing
else.

Throughout, $g\in G=\{s,p\}$ indexes nodes and directions (named by
destination), $n\in N$ products, $r\in R$ RTI types, $h\in H$ stack
types, $i,j\in I$ slot positions, $k\in K$ trucks and $t\in T$ periods,
with $\rho(n)$ the RTI type product $n$ travels in and
$M(r)=\{n:\rho(n)=r\}$; variables whose direction is fixed by
definition omit $g$. Table \ref{tab:sb-Nomenclature} is the
nomenclature.

\begin{longtable}[hbt!]{|p{0.15\textwidth}|p{0.79\textwidth}|}
    \caption{Nomenclature. The pricing problem of
    Section~\ref{sec:pricing} is one block of this model with the index
    $(g,k,t)$ deleted, so it reuses these symbols with that index
    dropped; the three variables it adds are listed last.}
    \label{tab:sb-Nomenclature}\\
    \hline
    \multicolumn{2}{|l|}{\textbf{Indices and sets}}\\\hline
    $G$            & Supplier and plant nodes, $G=\{s,p\}$; $g=s$ is the folded
                     direction (plant to supplier), $g=p$ the full one.\\
    $N$, $R$, $H$  & Products, RTI types, stack types.\\
    $I$            & Stack-slot indices on a truck floor, bounded by
                     \eqref{sb:eq:Ibound}.\\
    $K$, $T$       & Trucks available per direction and period; planning periods.\\
    $\rho(n)$, $M(r)$ & The RTI type product $n$ travels in; the products type $r$
                     carries, $M(r)=\{n:\rho(n)=r\}$.\\
    $R(h)$         & RTI types compatible with stack type $h$.\\
    \hline
    \multicolumn{2}{|l|}{\textbf{Parameters}}\\\hline
    $cK$           & Cost of dispatching one truck.\\
    $lt$, $SC$     & Lead time, $lt\ge1$; coverage window, in periods.\\
    $i_{gn}$, $ie_{gr}$ & Initial inventory at $g$ of full RTIs of $n$, of folded RTIs of $r$.\\
    $ic_{gn}$      & Inventory holding cost of $n$ at $g$.\\
    $d_{gnt}$      & Production plan for $n$ at $g$ in $t$: filled at $s$, consumed at $p$.\\
    $a_{nt}$, $ae_{rt}$ & Pipeline arrivals at $p$, respectively $s$, in $t\le lt$.\\
    $ss_{gnt}$, $sse_{gtr}$ & Coverage stock of full, respectively folded, RTIs
                     \eqref{sb:eq:coverage}.\\
    $w_{n}$, $w^{\mathrm{f}}_{r}$ & Weight of a full RTI carrying $n$; of a folded
                     RTI of type $r$. Both $>0$.\\
    $pK^{w}$, $W^{\max}$ & Truck payload; per-stack crush-limit cap on the
                     \emph{total} weight of one stack, over all its layers.\\
    $pK,qK,rK$     & Truck floor length, floor width, inner height;
                     $\Gamma_{K}=pK\cdot qK\cdot rK$.\\
    $\alpha$       & Minimum truck utilisation, on whichever of volume or payload
                     binds, $0<\alpha\le1$.\\
    $p_{h},q_{h}$  & Footprint length and width of stack type $h$.\\
    $r_{gr}$, $a_{r}$, $\gamma_{gr}$ & Height of one RTI of type $r$ at node $g$
                     (full at $p$, folded at $s$); its footprint area; its volume
                     $\gamma_{gr}=a_{r}r_{gr}$.\\
    $n_{hr}$       & RTIs of type $r\in R(h)$ per layer of stack type $h$; one layer
                     tiles the footprint, $n_{hr}a_{r}=p_{h}q_{h}$.\\
    $\hat w_{gr}$  & Weight per RTI in the layer-capacity form of the utilisation
                     floor: $\hat w_{sr}=w^{\mathrm{f}}_{r}$,
                     $\hat w_{pr}=\max_{n\in M(r)}w_{n}$.\\
    $\eta^{K},\eta^{I}$ & Weights of the transport and inventory objectives;
                     both $1$ throughout.\\
    $M^{x},M^{y}$  & Big-$M$ constants of the non-overlap rows, $M^{x}=pK$, $M^{y}=qK$.\\
    $\bar u^{V}_{n}$, $\bar u^{\mathrm{e}}_{r}$ & $\lfloor W^{\max}/w_{n}\rfloor$ full
                     RTIs of $n$ per slot; $\lfloor pK^{w}/w^{\mathrm{f}}_{r}\rfloor$
                     folded RTIs of $r$ per truck.\\
    \hline
    \multicolumn{2}{|l|}{\textbf{Decision variables --- the truck and its layout}}\\\hline
    $Z_{gkt}$        & $\in\{0,1\}$. Truck $k$ of direction $g$ in $t$ is configured,
                       i.e.\ carries a layout.\\
    $Y_{gkt}$        & $\in\{0,1\}$. That truck is dispatched.\\
    $\theta_{gkt}$   & $\in\{0,1\}$. Selects the resource on which its utilisation
                       floor is required: volume if $1$, payload if $0$.\\
    $X_{gihkt}$      & $\in\{0,1\}$. Slot $i$ holds a stack of type $h$.\\
    $L_{gihktr}$     & $\in\mathbb{Z}_{\ge 0}$. Layers of RTI type $r\in R(h)$ in slot $i$.\\
    $LX_{gihkt},LY_{gihkt}$ & $\in\{0,1\}$. Stack $h$'s length lies along the $x$-,
                       respectively $y$-axis in slot $i$.\\
    $PX_{gikt},PY_{gikt}$ & $\ge 0$. Bottom-left coordinates of slot $i$ on the floor.\\
    $A,B,C,D$ (on $gijkt$) & $\in\{0,1\}$. Which of the four separating
                       half-planes resolves the overlap of slots $i$ and $j$.\\
    \hline
    \multicolumn{2}{|l|}{\textbf{Decision variables --- load and inventory}}\\\hline
    $V_{kint}$       & $\in\mathbb{Z}_{\ge 0}$. Full RTIs of $n$ in slot $i$ of
                       plant-bound truck $k$ in $t$.\\
    $V^{\mathrm{e}}_{krt}$ & $\in\mathbb{Z}_{\ge 0}$. Folded RTIs of type $r$ on
                       supplier-bound truck $k$ in $t$.\\
    $I_{gnt}$, $Ie_{gtr}$ & $\in\mathbb{Z}_{\ge 0}$. End-of-period inventory of full
                       RTIs of $n$, of folded RTIs of $r$, at $g$.\\
    \hline
    \multicolumn{2}{|l|}{\textbf{Derived quantities} (abbreviations, \eqref{sb:eq:agg-V}--\eqref{sb:eq:agg-L})}\\\hline
    $\overline{V}_{nt}$, $\overline{V^{\mathrm{e}}}_{rt}$ & Full RTIs of $n$ delivered
                       to the plant in $t$; folded RTIs of $r$ returned to the supplier.\\
    $\overline{V}_{knt}$ & Full RTIs of $n$ on plant-bound truck $k$ in $t$.\\
    $\widehat{X}_{gikt}$, $\overline{L}_{ghktr}$ & Slot $i$ is used; layers of type $r$
                       under stack type $h$ on that truck, over all slots.\\
    \hline
    \multicolumn{2}{|l|}{\textbf{Added by the reformulation} (Sections \ref{sec:master}--\ref{sec:pricing})}\\\hline
    $\lambda^{g}_{\omega t}$ & $\in\mathbb{Z}_{\ge0}$. Trucks of class $(g,t)$
                       dispatched carrying column $\omega$.\\
    $\pi_{nt},\pi_{rt}$, $\mu_{gt}$ & Duals of \eqref{bp:MP-def-V}--\eqref{bp:MP-def-Ve},
                       respectively of \eqref{bp:MP-count} and of every further row a
                       dispatch variable of $(g,t)$ occupies.\\
    $n_{gt}$, $X$    & Dispatch total of class $(g,t)$, $n_{gt}=\sum_{\omega}
                       \lambda^{g}_{\omega t}$; the fleet $X=\sum_{g,t}n_{gt}$.\\
    \hline
\end{longtable}


\subsection{Derived quantities}

Five sums recur, as \emph{abbreviations} and not constraints.
\begin{align}
\overline{V}_{nt}&=\sum_{k\in K}\sum_{i\in I} V_{kint}
    &&\forall n\in N,\; t\in T
    \label{sb:eq:agg-V}\\
\overline{V^{\mathrm{e}}}_{rt}&=\sum_{k\in K} V^{\mathrm{e}}_{krt}
    &&\forall r\in R,\; t\in T
    \label{sb:eq:agg-Ve}\\
\overline{V}_{knt}&=\sum_{i\in I} V_{kint}
    &&\forall k\in K,\; n\in N,\; t\in T
    \label{sb:eq:agg-Vk}\\
\widehat{X}_{gikt}&=\sum_{h\in H} X_{gihkt}
    &&\forall g\in G,\; i\in I,\; k\in K,\; t\in T
    \label{sb:eq:agg-X}\\
\overline{L}_{ghktr}&=\sum_{i\in I} L_{gihktr}
    &&\forall g\in G,\; h\in H,\; k\in K,\; t\in T,\; r\in R(h)
    \label{sb:eq:agg-L}
\end{align}
The first two carry the method. The inventory block reads
$\overline{V}_{nt}$ and $\overline{V^{\mathrm{e}}}_{rt}$ and nothing
else about the trucks --- not which truck carried what, nor how the
load was placed.

\subsection{Objective}

Transport cost counts \emph{dispatches}, so a configured but
undispatched truck costs nothing; holding cost falls on erected RTIs.
\begin{align}
\min\; z^{K} &= \sum_{g\in G}\sum_{k\in K}\sum_{t\in T} cK\cdot Y_{gkt},
\qquad
\min\; z^{I} = \sum_{g\in G}\sum_{n\in N}\sum_{t\in T} ic_{gn}\cdot I_{gnt}
    \label{sb:eq:obj-zK}\\
\min\;& \eta^{K} z^{K}+\eta^{I} z^{I},
    \qquad \eta^{K}=\eta^{I}=1 \text{ throughout},
\label{sb:eq:obj}
\end{align}
\subsection{Inventory, flow balance and coverage}

Arrivals are delayed; inventory itself accumulates one period at a
time. For the plant,
\begin{align}
I_{pnt}  &=I_{p,n,t-1}-d_{pnt}+\overline{V}_{n,t-lt}
    &&\forall n\in N,\; t\ge lt+1
    \label{sb:eq:inv-plant-main}\\
I_{pnt}  &=I_{p,n,t-1}-d_{pnt}+a_{nt}
    &&\forall n\in N,\; t\in\{2,\dots,lt\}
    \label{sb:eq:inv-plant-transient}\\
I_{pnt}  &=i_{pn}-d_{pnt}+a_{nt}
    &&\forall n\in N,\; t=1
    \label{sb:eq:inv-plant-init}\\
Ie_{ptr} &=Ie_{p,t-1,r}+\sum_{n\in M(r)} d_{pnt}-\overline{V^{\mathrm{e}}}_{rt}
    &&\forall r\in R,\; t\ge 2
    \label{sb:eq:inve-plant}\\
Ie_{ptr} &=ie_{pr}+\sum_{n\in M(r)} d_{pnt}-\overline{V^{\mathrm{e}}}_{rt}
    &&\forall r\in R,\; t=1
    \label{sb:eq:inve-plant-init}
\end{align}
Consuming a full RTI produces a folded one in the same period, which is
why $d_{pnt}$ drives \eqref{sb:eq:inve-plant}. For the supplier,
\begin{align}
I_{snt}  &=I_{s,n,t-1}+d_{snt}-\overline{V}_{nt}
    &&\forall n\in N,\; t\ge 2
    \label{sb:eq:inv-sup-main}\\
I_{snt}  &=i_{sn}+d_{snt}-\overline{V}_{nt}
    &&\forall n\in N,\; t=1
    \label{sb:eq:inv-sup-init}\\
Ie_{str} &=Ie_{s,t-1,r}+\overline{V^{\mathrm{e}}}_{r,t-lt}-\sum_{n\in M(r)} d_{snt}
    &&\forall r\in R,\; t\ge lt+1
    \label{sb:eq:inve-sup-main}\\
Ie_{str} &=Ie_{s,t-1,r}+ae_{rt}-\sum_{n\in M(r)} d_{snt}
    &&\forall r\in R,\; t\in\{2,\dots,lt\}
    \label{sb:eq:inve-sup-transient}\\
Ie_{str} &=ie_{sr}+ae_{rt}-\sum_{n\in M(r)} d_{snt}
    &&\forall r\in R,\; t=1
    \label{sb:eq:inve-sup-init}
\end{align}
The lead time therefore acts on one side of a balance only --- the
sending node's stock falls in the dispatch period, the receiving node's
rises $lt$ periods later --- an alignment
Section~\ref{sec:branching} uses twice. Coverage stock is enforced at
both nodes and on both kinds of stock,
\begin{align}
I_{gnt}\ \ge\ ss_{gnt}
    &&\forall g\in G,\; n\in N,\; t\in T
    \label{sb:eq:safety-full}\\
Ie_{gtr}\ \ge\ sse_{gtr}
    &&\forall g\in G,\; r\in R,\; t\in T
    \label{sb:eq:safety-empty}
\end{align}
the requirement being that node's own series over the next $SC$
periods:
\begin{equation}
ss_{gnt}=\sum_{j=1}^{SC} d_{g,n,\tau(t,j)},
\qquad
sse_{gtr}=\sum_{n\in M(r)} ss_{gnt},
\qquad
\tau(t,j)=\begin{cases} t+j, & t+j\le |T|,\\ |T|-SC+j, & \text{otherwise.}\end{cases}
\label{sb:eq:coverage}
\end{equation}
Near the end of the horizon \eqref{sb:eq:coverage} wraps the window
back onto the last $SC$ periods rather than truncating it, following
\autocite{Diaz-Madronero2014APlanning}: truncating lowers the
requirement of the final periods and with it the optimal value, so the
tail rule belongs to the problem statement.

\subsection{The layout of a truck}

Everything from here to \eqref{sb:eq:stack-w-sup} is written for one
truck: one direction $g$, one index $k$, one period $t$. At most one
stack type occupies a slot and only a configured truck has slots; a
stack is oriented along one axis or the other; its height cannot exceed
the trailer's, which with $r_{gr}>0$ also empties an unused slot:
\begin{align}
  & \widehat{X}_{gikt}\;\le\;Z_{gkt} && \forall i\in I
    \label{sb:eq:slot-active}\\
  & LX_{gihkt}+LY_{gihkt}\;=\;X_{gihkt} && \forall h\in H,\; i\in I
    \label{sb:eq:orient}\\
  & \sum_{r\in R(h)} r_{gr}\,L_{gihktr}\;\le\;rK\cdot X_{gihkt}
    && \forall h\in H,\; i\in I
    \label{sb:eq:height}
\end{align}
Slots must not overlap --- for each pair $i<j$, the usual four-way
disjunction closed by a cover row --- and every slot lies inside the
floor:
\begin{align}
PX_{gikt}+\sum_{h\in H}\bigl(p_{h}LX_{gihkt}+q_{h}LY_{gihkt}\bigr)
    &\le PX_{gjkt}+M^{x}(1-A_{gijkt})
    \label{sb:eq:nonov-x1}\\
PX_{gjkt}+\sum_{h\in H}\bigl(p_{h}LX_{gjhkt}+q_{h}LY_{gjhkt}\bigr)
    &\le PX_{gikt}+M^{x}(1-B_{gijkt})
    \label{sb:eq:nonov-x2}\\
PY_{gikt}+\sum_{h\in H}\bigl(q_{h}LX_{gihkt}+p_{h}LY_{gihkt}\bigr)
    &\le PY_{gjkt}+M^{y}(1-C_{gijkt})
    \label{sb:eq:nonov-y1}\\
PY_{gjkt}+\sum_{h\in H}\bigl(q_{h}LX_{gjhkt}+p_{h}LY_{gjhkt}\bigr)
    &\le PY_{gikt}+M^{y}(1-D_{gijkt})
    \label{sb:eq:nonov-y2}\\
A_{gijkt}+B_{gijkt}+C_{gijkt}+D_{gijkt}
    &\;\ge\;\widehat{X}_{gikt}+\widehat{X}_{gjkt}-1
    \label{sb:eq:cover}\\
PX_{gikt}+\sum_{h\in H}\bigl(p_{h}LX_{gihkt}+q_{h}LY_{gihkt}\bigr)
    &\le pK
    \label{sb:eq:cont-x}\\
PY_{gikt}+\sum_{h\in H}\bigl(q_{h}LX_{gihkt}+p_{h}LY_{gihkt}\bigr)
    &\le qK
    \label{sb:eq:cont-y}
\end{align}
The containment rows are unconditional, which is what lets the
non-overlap rows take the tightened $M^{x}=pK$, $M^{y}=qK$.

\subsection{Full truckload on every dispatched truck}

A dispatched truck must be full, and it can be full in two ways: it may
run out of space or of payload. A volume-only floor is not a
conservative simplification but an infeasible one --- filling the
fraction $\alpha=0.55$ of a $97.92$~m$^{3}$ trailer
($13\,600\times2\,400\times3\,000$~mm) with freight at $579$~kg/m$^{3}$
weighs $31.2$~t against a $24$~t cap, and seven of the industrial
catalogue's fifteen types are denser still. Summing the two ratios is
the opposite error: it accepts a truck $35\%$ full by volume and $20\%$
by weight. We therefore require a fraction $\alpha$ of \emph{either}
resource, a binary $\theta_{gkt}$ selecting which one binds, on the
load the truck carries with each commodity's own weight:
\begin{align}
\sum_{n\in N}\gamma_{p\rho(n)}\,\overline{V}_{knt}
&\;\ge\;\alpha\,\Gamma_{K}\bigl(Y_{pkt}-1+\theta_{pkt}\bigr)
    &&\forall k\in K,\; t\in T
    \label{sb:eq:util-p}\\
\sum_{n\in N} w_{n}\,\overline{V}_{knt}
&\;\ge\;\alpha\,pK^{w}\bigl(Y_{pkt}-\theta_{pkt}\bigr)
    &&\forall k\in K,\; t\in T
    \label{sb:eq:utilw-p}\\
\sum_{r\in R}\gamma_{sr}\,V^{\mathrm{e}}_{krt}
&\;\ge\;\alpha\,\Gamma_{K}\bigl(Y_{skt}-1+\theta_{skt}\bigr)
    &&\forall k\in K,\; t\in T
    \label{sb:eq:util-s}\\
\sum_{r\in R} w^{\mathrm{f}}_{r}\,V^{\mathrm{e}}_{krt}
&\;\ge\;\alpha\,pK^{w}\bigl(Y_{skt}-\theta_{skt}\bigr)
    &&\forall k\in K,\; t\in T
    \label{sb:eq:utilw-s}
\end{align}
Exactly one of each pair is active at $Y_{gkt}=1$, so each is the
disjunction ``full by volume \emph{or} full by payload''. The same
disjunction on the layer \emph{capacity} the stacks provide rather than
on the load they receive,
\begin{align}
  \sum_{h\in H}\sum_{r\in R(h)}
    \gamma_{gr}\,n_{hr}\,\overline{L}_{ghktr}
&\;\ge\;\alpha\,\Gamma_{K}\bigl(Y_{gkt}-1+\theta_{gkt}\bigr),
\quad
  \sum_{h\in H}\sum_{r\in R(h)}
    \hat w_{gr}\,n_{hr}\,\overline{L}_{ghktr}
\;\ge\;\alpha\,pK^{w}\bigl(Y_{gkt}-\theta_{gkt}\bigr)
    \label{sb:eq:util-cap}
\end{align}
is a valid tightening, kept because it lets a solver establish a layout
before choosing a product assignment. Which form is the
\emph{requirement} matters.
\eqref{sb:eq:util-p}--\eqref{sb:eq:utilw-s} is, full truckload being a
property of the freight that moves; \eqref{sb:eq:util-cap} follows from
it because a slot never receives more RTIs of type $r$ than its layers
hold and $\hat w_{gr}\ge w_{n}$ for every $n\in M(r)$, and taken
\emph{alone} would admit trucks that are not full. That inequality is
also why $\hat w_{pr}=\max_{n\in M(r)}w_{n}$ is the family
\emph{maximum}: the family \emph{minimum} would assert that
\emph{every} product assignment of this layout fills the truck, cutting
off feasible loads.

\subsection{Dispatch, loading and weight}

A dispatched truck must be configured,
\begin{equation}
Y_{gkt}\;\le\;Z_{gkt}\qquad\forall g\in G,\;k\in K,\;t\in T,
\label{sb:eq:lc-active}
\end{equation}
while the converse fails: $Z_{gkt}=1$ with $Y_{gkt}=0$ is a truck
carrying a layout and no freight, which in the reformulated block is
exactly the empty column (Section~\ref{sec:pricing}).

Full RTIs in a slot cannot exceed the layer capacity of the stack
placed there, products sharing an RTI type partitioning it; folded RTIs
are capped by the truck's aggregate capacity, folded weight being
deterministic in the type. Both loads are gated on the dispatch, at
their own analytic bounds rather than a generic big-$M$:
\begin{align}
  & \sum_{n\in M(r)} V_{kint}
      \;\le\;\sum_{h\,:\,r\in R(h)} n_{hr}\,L_{p,i,h,k,t,r}
    && \forall k\in K,\; i\in I,\; r\in R,\; t\in T
    \label{sb:eq:slot-cap}\\
  & V^{\mathrm{e}}_{krt}
      \;\le\;\sum_{h\,:\,r\in R(h)} n_{hr}\,\overline{L}_{s,h,k,t,r}
    && \forall k\in K,\; r\in R,\; t\in T
    \label{sb:eq:empty-cap}\\
  & V_{kint}\le \bar u^{V}_{n}\,Y_{pkt}
    && \forall k\in K,\; i\in I,\; n\in N,\; t\in T
    \label{fix:slot-link}\\
  & V^{\mathrm{e}}_{krt}\le \bar u^{\mathrm{e}}_{r}\,Y_{skt}
    && \forall k\in K,\; r\in R,\; t\in T
    \label{fix:empty-link}
\end{align}
The payload is a property of the truck, the crush limit of each stack;
on plant-bound trucks that cap is \emph{period-exact}, evaluated on the
mix actually loaded in the slot, on supplier-bound trucks deterministic
in the layout:
\begin{align}
  & \sum_{n\in N} w_{n}\,\overline{V}_{knt}\;\le\;pK^{w}\,Y_{pkt},
  \qquad
    \sum_{r\in R} w^{\mathrm{f}}_{r}\,V^{\mathrm{e}}_{krt}
      \;\le\;pK^{w}\,Y_{skt}
    && \forall k\in K,\; t\in T
    \label{sb:eq:truck-w-plant}\\
  & \sum_{n\in N} w_{n}\,V_{kint}\;\le\;W^{\max}
    && \forall k\in K,\; i\in I,\; t\in T
    \label{sb:eq:stack-w-plant}\\
  & \sum_{h\in H}\sum_{r\in R(h)}
      w^{\mathrm{f}}_{r}\,n_{hr}\,L_{s,i,h,k,t,r}\;\le\;W^{\max}
    && \forall i\in I,\; k\in K,\; t\in T
    \label{sb:eq:stack-w-sup}
\end{align}
where \eqref{sb:eq:truck-w-plant} names both payload rows and
\eqref{fix:slot-link} is tighter in the relaxation than their
implication, since what bounds a slot's content is the crush cap.

\subsection{Symmetry, slot-set sizing and data assumptions}
\label{sb:sec:symm}

The slots of a truck are interchangeable, and so are the trucks of a
direction and period. Both are ordered:
\begin{align}
  & \widehat{X}_{g,i+1,k,t}\;\le\;\widehat{X}_{gikt}
    && \forall g\in G,\; k\in K,\; t\in T,\; i<|I|
    \label{sb:eq:sym}\\
  & Y_{g,k+1,t}\;\le\;Y_{gkt}
    && \forall g\in G,\; t\in T,\; k<|K|
    \label{dec:sym-truck}
\end{align}
Re-indexing makes \eqref{dec:sym-truck} valid, but it removes only
\emph{which} subset of trucks is used: loads and layouts stay
interchangeable among those used, so a plan dispatching $m$ trucks in a
period still appears $m!$ times --- the difficulty
Section~\ref{sec:decomp} removes by aggregation rather than by a row.
The slot count is a modelling parameter, and
\begin{equation}
|I|\;\le\;\Bigl\lfloor
    (pK\cdot qK)\big/\min_{h\in H} p_{h}q_{h}
\Bigr\rfloor
\label{sb:eq:Ibound}
\end{equation}
eliminates provably unused slot variables without cutting any feasible
layout. It is more than preprocessing: a placement model given fewer
slots than a layout needs can report infeasible while the load is
placeable (Section~\ref{sec:branching}).

Four properties of the data are assumed: $\rho$ is a partition;
$lt\ge1$, which the lagged terms require; all weights and heights are
strictly positive, which makes $\bar u^{V}_{n}$ and
$\bar u^{\mathrm{e}}_{r}$ well defined and lets \eqref{sb:eq:height}
empty an unused slot; and the footprint geometry is consistent, one
layer of stack type $h$ tiling its footprint, $n_{hr}a_{r}=p_{h}q_{h}$,
with the volume charged to an RTI its bounding box,
$\gamma_{gr}=a_{r}r_{gr}$. \emph{The last is a departure}, and
Proposition~\ref{prop:inclusion} and the floor-area inequality
\eqref{bp:PP-areacut} are both false without it. It costs nothing
here, both instance families building one stack type per distinct RTI
footprint, so $n_{hr}=1$ and the relation is an identity.

\subsection{Domains}

\begin{align*}
&Z_{gkt},\,Y_{gkt},\,\theta_{gkt},\,X_{gihkt},\,LX_{gihkt},\,LY_{gihkt},\,
  A_{gijkt},\,B_{gijkt},\,C_{gijkt},\,D_{gijkt}\in\{0,1\},\\
&L_{gihktr},\,V_{kint},\,V^{\mathrm{e}}_{krt},\,I_{gnt},\,Ie_{gtr}
  \in\mathbb{Z}_{\ge 0},\qquad
PX_{gikt},\,PY_{gikt}\ge 0.
\end{align*}
